% 缺口 4: 架构组件消融
\begin{table}[H]
\centering
\caption{CNN-Transformer架构组件消融：各组件的独立与联合贡献（full特征集，跨2种子）}
\label{tab:arch_ablation}
\small
\setlength{\tabcolsep}{3pt}
\begin{tabular}{lccc}
\toprule
架构变体 & 参数量 & Sharpe & IC \\
\midrule
CNN-Transformer（完整）& 1.75M & $\bf 5.46$ & $\bf 0.152$ \\
CNN-only（去Transformer）& 1.02M & $0.04$ & $0.014$ \\
Transformer-only（去CNN）& 0.72M & $4.30$ & $0.143$ \\
Transformer-only（参数量匹配，1.81M）& 1.81M & $5.08$ & $0.142$ \\
\bottomrule
\multicolumn{4}{p{\textwidth}}{\footnotesize 注：跨2种子（0,1）均值。CNN-only完全丧失预测能力（Sharpe$\approx$0），表明纯卷积的局部感受野不足以捕获跨源交互。Transformer-only即使参数量匹配（1.81M vs.\ 1.75M），Sharpe仍低于完整混合架构（$5.08$ vs.\ $5.46$），证实"混合"（CNN局部提取$\to$Transformer全局交互）本身是增益的本质来源，参数量差异非混淆因素。}\\
\end{tabular}
\end{table}
